\chapter{相关技术简介}

本章将介绍实现本系统所涉及的关键技术。为了支撑AI多任务场景下的复杂调度需求，本系统在选型上重点考虑了分布式处理能力、异构数据适配性以及系统的高可扩展性。

\section{基础开发与运行环境}

\subsection{Python编程语言}

Python是用可移植的ANSIC写成的，可以在今天所用的各主要平台上运行。尽管是面向对象的程序设计语言，程序员也可以根据需要，利用Python，选择进行面向对象的程序设计或进行结构化的程序设计\cite{11, 13}。Python的创始人为Guido van Rossum。1989年圣诞节期间，在阿姆斯特丹，Guido为了打发圣诞节的无趣，决心开发一个新的脚本解释程序，做为ABC语言的一种继承\cite{12}。Python 是用可移植的ANSIC写成的，可以在今天所用的各主要平台上运行。

\subsection{Docker与Kubernetes技术}

Docker是一种开源的容器化平台，能够将应用及其依赖打包为轻量级、可移植的容器，并在任意支持Docker的环境中保持一致运行，从而实现应用的快速部署与分发\cite{8}。通过将各服务封装于Docker容器，可简化部署、扩展与管理流程，保障系统在不同环境下的稳定运行。在此基础上，容器化技术还能针对不同形态与领域的云服务及云平台产品，统一其API、功能组件、业务服务与操作流程，为用户提供贯通一致、互联互通、便捷快速的云平台使用体验，并提升云边协同的容器云管控能力。但在实际的容器化部署过程中，容器编排与资源管理仍是颇具挑战的环节。以本系统为例，当鉴权服务器、调度服务器等多个服务被容器化后，如何合理编排这些容器、实现高效的资源分配与负载均衡，从而保障系统在不同负载状况下均能稳定运行，成为亟待解决的问题。

Docker镜像是Docker容器的基础，它是一个只读的模板，包含了运行应用程序所需的一切，如代码、运行时环境、库、工具等。Docker镜像的打包高度依赖Dockerfile配置文件，这个文件包含了一系列的指令，用于描述如何构建 Docker 镜像。本系统迭代过程中，每个测试版和稳定版程序都会被打包成镜像并发布到对应的环境，其打包过程如下：首先要明确镜像内容，将本次迭代所依赖的基础环境和pip软件包写进配置文件中，接着编写或修改Dockerfile，若本次迭代有新增目录，需相应修改以将这些新增部分加入镜像，同时编写或修改.dockerignore文件，指明仅供开发需要、可在运行环境中剔除的文件以减小镜像体积和复杂度，之后通过持续集成系统自动运行命令构建镜像，再在测试环境中进行持续集成自动测试和人工检测，当前置环节都通过后，该镜像即满足本次迭代的准出要求，可通过腾讯公司的持续部署系统轻松部署到生产环境，Docker镜像大大便利了本系统的迭代与上线工作。

此外，还有Dev Container技术。Dev Container即开发容器，是一种基于容器技术的开发环境解决方案。它借助Docker等容器化技术，将开发所需的运行时环境、工具、依赖等进行打包，形成一个独立且可移植的开发环境。开发人员可以在这个容器内进行代码编写、测试和调试等工作，而无需担心本地环境与项目环境不兼容的问题。本系统采用Dev Container作为开发环境，开发人员不用特意配置环境，可以大大缩短项目交接时间。

Kubernetes简称k8s，是一个开源的容器编排系统，用于自动化部署、扩展和管理容器化应用\cite{8}。腾讯公司内部的持续部署系统基于k8s，有着高效的容器编排机制，可以方便地部署本系统。近年来，微服务架构(Microservices Architecture)受到了业界和学术界的广泛关注，成为软件系统开发的另一个趋势，有着较为广阔的应用前景\cite{7}，而k8s与Docker配合，可以快速地部署大部分微服务程序。

\section{分布式任务调度架构}

\subsection{Celery分布式任务队列}

Celery作为高性能分布式任务队列框架，其核心通过异步消息机制实现任务调度与执行的解耦。针对AI多任务场景下任务量大、计算密集且耗时较长的特性，Celery利用中间人（Broker）协调客户端与消费者（Worker），不仅提供任务调度、重试及状态监控等基础功能，更具备卓越的横向扩展能力，每分钟可处理数以百万计的任务。由于AI大模型相关算法多基于Python开发，选用原生支持Python生态的Celery可深度适配复杂的参数类型与调用约定，有效解决了第一子问题中提到的系统适配性不足及数据处理能力差的痛点。

在本系统中，Celery架构由生产者、消息队列与消费者共同构成，完美契合AI场景下多对多的任务交互特性。其中，“API接入层”与“任务调度层”扮演生产者的角色，负责解析业务逻辑并封装任务参数；“任务执行层”则作为分布式消费者，负责具体的算力消耗与结果上报。通过采用异步协程gevent执行池配合Redis消息队列，系统能够对AI任务进行高效并发调度，实现了由串行执行向并行处理的转变，直接解决了第三子问题中因资源互斥导致的计算资源浪费难题，显著提升了整体执行效率\cite{14}。

\subsection{Redis消息中间件}

Redis使用C语言编写，作为一种基于内存的非关系型数据库，以其高性能、灵活的数据结构、可扩展性和数据持久化机制等特性，成为管理分布式存储系统元数据的理想选择\cite{16, yin2025}。Redis核心引擎是单线程的，所以不存在多线程数据冲突、竞争的问题，十分适合本系统多对多、高并发的场景，因此本系统选用Redis作为消息中间件。

一个完整的Celery工程需要两个消息队列实例，一个负责存储待执行的任务，另一个负责存储任务执行结果。本系统选用同一个Redis实例的两个数据库（db-number）分别存储待执行任务和执行结果。从后续实践来看，这样的设计足以承载本系统的任务压力。

\section{接口服务与远程通信技术}

\subsection{FastAPI服务框架}

FastAPI是一个用于构建Web API的Python框架，适合开发高性能的RESTful API。FastAPI后端主要用于接收前端传递的数据并选取合适的函数、向前端返回最后的计算结果\cite{15}。通过结合静态类型检查和异步编程，以及利用现代Python的语言特性，FastAPI提供了一个高性能、易于使用和可维护的Web框架。本系统的上下游系统均以Python编写，因此采用Python生态体系内的Web框架能更好地适配整个工程体系，兼容各种参数类型和调用习惯。

同时，项目组选定Vue.js作为本系统的设计实现框架。响应式数据是Vue中MVVM模式的实现，也就是ViewModel，它可以对应多个视图，具有可重用性，提高了数据的绑定使用效率；同时，在本系统设计中，它的虚拟DOM渲染设计，减少了对真实DOM的频繁操作，提升了页面渲染效率\cite{20}。

\subsection{tRPC过程调用技术}

远程过程调用（RPC）允许客户计算机通过网络调用服务器计算机程序或方法，且如同在本地执行般透明，极大简化了网络分布式系统的编程工作，广泛应用于云计算、大数据、微服务等领域，帮助开发者便捷地构建高性能、可扩展的分布式应用程序\cite{2}。

tRPC是腾讯自研的高性能、跨平台、插件化、具备高度服务治理能力的RPC框架。本系统服务腾讯公司的AI生态，为了和上下游保持一致，也采用了tRPC作为系统间RPC的标准协议。tRPC具有丰富的特性，首先它具备跨语言能力，因使用跨语言的序列化与反序列化库Protocol Buffers作为序列化组件；其次支持多通信协议，除默认的tRPC协议外，还兼容Thrift、gRPC等协议，能为异构系统构建提供技术支持；同时支持流式RPC，适用于大文件上传/下载、消息推送、AI类语音识别/视频理解等场景，而本系统需承接大量AI数据传输与交互需求，因此该特性尤为重要；此外，它拥有丰富的插件生态，不仅提供大量微服务相关组件以帮助开发者快速构建治理体系，还依托腾讯公司背景包含内部生态治理组件，在本系统构建中发挥不可替代的作用；并且具备良好的可扩展性，基于插件化设计，用户可通过二次开发扩展请求参数校验、鉴权、请求录制等框架能力；最后，它还提供流控和过载保护功能，通过多种流量控制与过载保护插件防止服务因访问突增而无法使用，为系统稳定性提供了有力支持。

\section{数据持久化与版本管理}

\subsection{PostgreSQL数据库}

本系统选用PostgreSQL作为数据持久化存储系统。PostgreSQL是由加州大学伯克利分校计算机科学系开发的对象关系数据库管理系统\cite{3}。相较于MySQL，它在大数据处理方面具有显著优势，具备更强大的功能、更丰富的索引类型、更高效的并行查询能力以及更灵活的物化视图等特性\cite{18}。PostgreSQL支持大部分SQL标准，并提供众多现代功能，如复杂查询、外键、触发器、可更新视图、交易完整性以及多版本并发控制等。用户还能通过多种途径对其进行扩展，像添加新的数据类型、函数、运算符、集合函数、索引方法以及程序语言等\cite{4}。在数据安全性与灵活性方面，PostgreSQL 通过多版本并发控制实现快照隔离 \cite{17}，保障数据安全。同时，它具备在运行时添加新的数据类型、函数和访问方法的机制\cite{19}，且能确保事务安全\cite{9520330}。此外，PostgreSQL的数据库表结构可更新，还支持JSON作为数据类型。PostgreSQL展现出的诸多特性都很适配本系统，它能充分满足AI任务所面临的海量数据处理需求，与AI多样化的数据处理要求高度契合，多版本并发控制机制也能有力地保障了数据安全。 

\subsection{SQLAlchemy与Alembic}

SQLAlchemy是Python编程语言下的一款开源软件。 提供了SQL工具包及对象关系映射工具，使用MIT许可证发行\cite{5}。SQLAlchemy具有多方面的优势，它包含ORM组件，能够将Python的数据结构自动映射为数据库表，实现了数据库与数据结构的分离，这不仅屏蔽了数据库的底层差异，还为开发者提供了一套统一的API，便于开发与迁移；同时，它拥有强大灵活的API，结合Python自身丰富的语法，使得编写查询语句的工作变得十分简单，并且还支持众多数据库。

Alembic是一个轻量级的数据库迁移工具，由SQLAlchemy项目团队开发，主要用于管理SQLAlchemy项目中的数据库模式变更。它提供了一种版本控制的方式来管理数据库架构的演变，使得开发者可以方便地在不同版本的数据库模式之间进行迁移，确保数据库结构与应用程序代码的一致性。Alembic库具有多方面特点，它采用类似Git History的版本管理机制，为每个数据库版本编号并记录前后版本关系和版本间的迁移脚本，这一完善的版本控制机制不仅方便数据库自动更新，也使得在出故障时可以很方便地手动回滚表结构；同时，Alembic支持所有常见的数据库类型，还能够自动识别工程中数据结构的变更，并自动生成迁移脚本和版本信息。本系统采用Alembic作为数据库版本管理工具，在开发与迭代工作中，回滚了数个设计不当的数据结构，提高了开发效率与容错率。

\section{本章小结}

本章详细介绍了支撑系统运行的各项技术。上述技术并非孤立存在，而是通过有机整合，共同服务于解决AI任务调度的适配性、自动化及资源互斥等核心难题。Python与FastAPI负责接入层的逻辑校验，Celery与Redis构成了调度层的自动决策底座，tRPC保障了跨节点的可靠传输，而PostgreSQL则为资源互斥检测提供了精准的数据支撑。这些技术的应用为后续章节中系统的架构设计与详细实现提供了坚实保障。

\makeatletter

\@openrightfalse

\makeatother